\clearpage
\fancyhead[L]{\small FLORA.}
\SourceLeaf{481}
\begin{center}
\rule{0.55\linewidth}{0.4pt}\\[2pt]
Astronomische Nachrichten. Band XXVI. Nr. 622. S. 352. 10. Januar 1848.\\[2pt]
\rule{0.55\linewidth}{0.4pt}
\end{center}

Von Geheimrat \textsc{Gauss} habe ich die folgenden Beobachtungen der \emph{Flora} am Göttinger Meridiankreis erhalten.

\begin{center}
\setlength{\tabcolsep}{4pt}
\begin{tabular}{l|r r r|r|r}
Dec. 10 & $11^{\rm u}$ & $11^{\rm m}$ & $4^{\rm s}0$ & $66^\circ\ 47'\ 26''2$ & $+14^\circ\ 22'\ 50''1$ \\
\hfill 11 & & $6$ & $7.9$ & $66\ 32\ 5.4$ & $14\ 25\ 42.6$ \\
\hfill 12 & & $1$ & $11.0$ & $66\ 16\ 59.9$ & $14\ 28\ 40.0$ \\
\end{tabular}
\end{center}

Geheimrat \textsc{Gauss} bemerkt, dass die Rektaszension am 10. etwas weniger zuverlässig sei als die am 11. und dass Flora viel heller als Iris erscheine. Die Beobachtungen konnten bei voller Beleuchtung gemacht werden, und der Planet bietet über $\rho$ Tauri, der ihm einige Minuten weiter nördlich um $7^{\rm s}$ folgte, einen schönen Anblick.

